% ==============================================================================
% COLORS: BẢNG MÀU DÙNG CHUNG
% ------------------------------------------------------------------------------
%  Mỗi màu có 3 sắc độ:   <tên>  (chuẩn)   <tên>dam (đậm)   <tên>nhat (nền nhạt)
%  Ví dụ: \textcolor{xanhla}{...}, \colorbox{camnhat}{...}, colback=tim nhat
% ==============================================================================

% ---------- Màu thương hiệu Bách Khoa ----------
\definecolor{bkxanh}{HTML}{0B3D91}       % xanh navy chủ đạo (tiêu đề, khung)
\definecolor{bkxanhsang}{HTML}{1E88E5}   % xanh logo (điểm nhấn, liên kết)
\definecolor{bkxanhnen}{HTML}{EAF2FC}    % nền xanh rất nhạt (bảng sọc, khung)
\definecolor{bkvien}{HTML}{0000B4}       % xanh khung bìa
\definecolor{bkxam}{HTML}{5F6B7A}        % chữ phụ

% ---------- Màu cơ bản (tên không dấu) ----------
\definecolor{xanhduong}{HTML}{1565C0}
\definecolor{xanhla}{HTML}{2E7D32}
\definecolor{ngoc}{HTML}{00838F}          % xanh ngọc / teal
\definecolor{do}{HTML}{C62828}
\definecolor{cam}{HTML}{E65100}
\definecolor{vang}{HTML}{F9A825}
\definecolor{tim}{HTML}{6A1B9A}
\definecolor{hong}{HTML}{AD1457}
\definecolor{nau}{HTML}{6D4C41}
\definecolor{xam}{HTML}{616161}

% ---------- Tự sinh sắc độ đậm / nhạt ----------
\newcommand{\TaoSacDo}[1]{%
  \colorlet{#1dam}{#1!75!black}%
  \colorlet{#1nhat}{#1!9!white}}
\TaoSacDo{xanhduong}\TaoSacDo{xanhla}\TaoSacDo{ngoc}\TaoSacDo{do}
\TaoSacDo{cam}\TaoSacDo{vang}\TaoSacDo{tim}\TaoSacDo{hong}
\TaoSacDo{nau}\TaoSacDo{xam}\TaoSacDo{bkxanh}
\colorlet{vangnhat}{vang!22!white}        % nền highlight vàng rõ hơn

% ---------- Màu cho mã nguồn ----------
\definecolor{manen}{HTML}{F7F9FC}
\definecolor{matukhoa}{HTML}{1D4ED8}
\definecolor{machuthich}{HTML}{15803D}
\definecolor{machuoi}{HTML}{9333EA}
\definecolor{maso}{HTML}{94A3B8}

% ---------- Lệnh tô chữ nhanh ----------
\newcommand{\chuxanh}[1]{\textcolor{xanhduong}{#1}}
\newcommand{\chula}[1]{\textcolor{xanhla}{#1}}
\newcommand{\chudo}[1]{\textcolor{do}{#1}}
\newcommand{\chucam}[1]{\textcolor{cam}{#1}}
\newcommand{\chutim}[1]{\textcolor{tim}{#1}}
\newcommand{\chuxam}[1]{\textcolor{xam}{#1}}
